% Modernised reading — fonds Alexandre Grothendieck, shelfmark 73, pages 1-28
% (the whole folder).
%
% This is an INTERPRETATION, not a transcription. It re-reads the pages in
% today's notation and vocabulary, states the hypotheses the manuscript leaves
% implicit, and drops the critical apparatus so that the mathematics reads
% straight through. Everything the brackets and underlines carried moves into
% a footnote. It is held to being mathematically correct as it stands; where
% the manuscript is loose or mistaken, the true statement is what appears here
% and the footnote says what the page has. In any dispute of substance the
% transcription governs: whoever cites Grothendieck cites batch-01.fr.tex
% (pp. 1-20) or batch-02.fr.tex (pp. 21-28).
%
% Pass: Opus 5.5 (claude-opus-5-5), 2026-10-10 — first pass, unchecked against the pages by a human.
% The folder was read whole, both batches together; this is one document for
% the shelfmark. It was made on Opus 5.5 and has not been measured against the
% reference reading of folder 115 (made on Opus 5). The literature named in the
% footnotes (Gauss's pentagon formula, Lagrange multipliers) is cited from
% memory and was not checked against the sources. Every identity the reading
% asserts beyond the page (Gauss's equation, the corrected discriminant, the
% P, Q, R recurrence and its range, the symmetric-square identity, the
% dimensions of N(u_*) for closed n-gons, n = 3..7, the collinear
% counter-example, the multiplier of the nine-gon of p. 18) was checked by a
% symbolic computation outside the file.
%
% Scope. Pages 2, 7, 11 and 13 are blank (the transcription skips them, and
% the blank pages were checked on the facsimile). The inventory's « lettre
% (s.d.) » is on no scanned leaf; nothing of it is described. Page 8 carries a
% short non-mathematical list in the margin, left out by the transcription and
% here.
%
% Spine. The map that sends a plane n-gon to the signed areas of its n « ears »
% (the triangles s_{i-1} s_i s_{i+1}), and its critical points:
%   p. 1        the functor Polg_n of n-gons in an affine plane (vertices and
%               side-lines), degenerations;
%   pp. 3-4     the multiplier relation in its first, « Lagrange » form
%               lambda_i u_{i-1} - lambda_{i+1} u_{i+1} = w (constant), and its
%               differences M_i (the four-term relation of pp. 20-28);
%   pp. 5-6     quadrilaterals: one linear relation among the ear areas, and
%               M_4 = P^2 with the four degenerate quadrilaterals as a frame;
%   pp. 8-10    an unrelated piece, written backwards: configurations,
%               operations, residual sets, tests, « cribles » (adaptive
%               questioning strategies);
%   p. 12       relations among lambda, mu, alpha for a closed n-gon;
%   pp. 14-15   the Lagrange formulation written out (sum lambda_i df_i =
%               mu sum dx_i + nu sum dy_i), and the closed formula, boxed;
%   pp. 16-17   quadrilaterals again: lambda_0 + lambda_1 = 0, an explicit
%               chart, the two diagonals;
%   p. 18       lambda_0 = 0: a star nine-gon;
%   p. 19       pentagons: a quadratic equation, which is Gauss's;
%   pp. 20-22   the four-term relation, u_i in the basis u_0, u_1, u_2, the
%               recurrence for P_i, Q_i, R_i, the closing equations (0)-(3);
%   pp. 23-27   the degenerate case, the Proposition and its proof;
%   p. 28       the Theorem: the four cases for N(u_*).
%
% Notation fixed for the whole document.
%  - k a field, V a vector plane over k, E an affine plane of direction V;
%    vertices s_i, sides u_i = s_{i+1} - s_i, i in Z/nZ; u_{ab} = u_a + ... +
%    u_{b-1} = s_b - s_a (the page's own notation, p. 15).
%  - mu_i = u_{i-1} ^ u_i in Lambda^2 V, identified with k by a choice of
%    volume form: twice the signed area of the ear at s_i. The pages write
%    mu_i on pp. 5 and 12 and f_i on pp. 14-19; mu_i throughout here, since
%    the page's f_i is in the same pages the name of nothing else. phi the map
%    u_* |-> (mu_i)_i.
%  - w for the constant vector of the Lagrange form (the page's u on pp. 3, 4
%    and 12, (nu, -mu) in coordinates on p. 14), so that u stays free.
%  - (L) and (*) for the two forms of the multiplier relation; N(u_*) the
%    space of multipliers (the page's, p. 25). (*) is used with the indexing of
%    pp. 25-28; p. 20 writes the same family shifted by one.
%  - The juxtaposition u v of two vectors on pp. 3, 4, 15, 18 and 28 is read
%    as the product in the symmetric algebra Sym(V) and written u.v. The
%    transcription's header likens it to the determinant; that reading is
%    shown false (p. 15, i = 2: the sign of the lambda_0 term) in a footnote.
%
% Substantive departures from the pages, each footnoted where it occurs:
%  - p. 16: the first displayed result for the vertex s_2 has the wrong sign;
%    the vertex is s_2, not the page's s_1. The boxed lambda_0 + lambda_1 = 0
%    stands.
%  - p. 19: the last expansion of the discriminant reads + 4 f_2 f_3; it is
%    - 4 mu_2 mu_3. The middle line (read with doubt) is not an identity.
%  - pp. 21-23: the coefficient polynomials of the closure relation (0) lie in
%    Z[Lambda_1, ..., Lambda_{n-1}], not Z[Lambda_1, ..., Lambda_{n-2}]; the
%    third closing equation (3) is written with the minus sign the page leaves
%    as « (-) ».
%  - p. 24: the indices of u_7, u_8, u_9 do not follow the rule written above
%    them; the argument is restated, not the list.
%  - pp. 25-27 (Proposition): the hypothesis that the u_i do not all lie on one
%    line (the margin of p. 23, « OPS u_0, u_1 lin. indép. ») is supplied: the
%    segment run back and forth, (e, -e, e, -e), has dim N = 2 with all
%    multipliers non-zero and satisfies (b) but not (c), (d). n = 0 mod 3 is
%    put into (d), where the page struck it. The injectivity of N(u_*) -> k^3
%    is replaced by that of N(u_*) -> k^2 on two consecutive components, which
%    needs only u_i != 0 and no hypothesis on k; the page's route is footnoted.
%    The proof of d) => a), announced and not written, is supplied, and needs
%    card k >= 3.
%  - p. 28 (Theorem): the equivalence (b_2) <=> (d) needs the same
%    non-collinearity; the exhaustiveness of the cases needs k infinite (the
%    « Lemme » of the margin of p. 26). The half-illegible gloss of (a) is
%    reconstructed as « phi is a submersion at u_* »; the criterion for (c)
%    (« il faut et il suffit », read with doubt) is proved here and holds.
%
% Modern names given and footnoted as ours: Lagrange multipliers and critical
% points, conormal space, symmetric square, projective frame, Gauss's pentagon
% formula, decision trees. None is on the pages; nothing on them cites anyone.
%
% Announced and never established: the open U of p. 1 (its condition is
% struck and unreadable); the dimension counts of p. 4 (prose lost); « épi ? »
% (p. 5), answered here from the formulas of the same page; the system (0)-(3)
% of p. 22, never solved; the meaning of CI_n^* on p. 12; the purpose of the
% quadric at the top of p. 20; the proof of d) => a) (p. 27); the marginal note
% of p. 28, largely lost. The « cribles » of pp. 8-10 are not applied to
% anything in the folder.
\documentclass[11pt,a4paper]{article}
\input{../preamble/grothendieck}

\folder{73}
\pages{1}{28}
\dating{[à partir de 1976]}
\watermark{Édition de démonstration\\interprétation personnelle de l'œuvre}
\foldertitle{Polygones : notes manuscrites (s.d.), lettre (s.d.) — lecture modernisée du dossier entier}

\begin{document}

\section*{Résumé}

\begin{resume}
Ces feuillets parlent de polygones du plan, et d'une question qui paraît
d'abord naïve : que sait-on d'un polygone quand on connaît les aires de ses
« oreilles » ? À chaque sommet on associe le triangle qu'il forme avec ses deux
voisins. Si l'on déforme le plan par une transformation affine --- une
application qui garde les droites et le parallélisme, mais ni les longueurs ni
les angles ---, toutes ces aires sont multipliées par un même nombre. Leurs
rapports sont donc des invariants du polygone en géométrie affine, où il n'y a
plus de distances mais où les rapports d'aires gardent un sens. Combien
d'information portent-ils, et où en perdent-ils ?

On peut penser aux polygones articulés, faits de tiges rigides reliées par des
charnières : on connaît les longueurs des côtés, et l'on se demande quelles
formes le polygone peut prendre. Ici les longueurs sont remplacées par les
aires des oreilles, et la géométrie euclidienne par la géométrie affine.

Pour le quadrilatère la réponse est complète. Les quatre aires vérifient une
seule relation --- l'aire du quadrilatère se lit aussi bien sur les deux
oreilles opposées d'une diagonale que sur celles de l'autre ---, et à cette
relation près elles déterminent le quadrilatère : les quadrilatères, à
transformation affine près, forment exactement un plan projectif, dans lequel
les quatre quadrilatères dégénérés, ceux dont un côté est réduit à un point,
sont quatre points en position générale. Pour le pentagone, le calcul aboutit à
une équation du second degré : deux pentagones ont en général les mêmes
oreilles. Cette équation est, sous un déguisement, une formule ancienne qui
relie l'aire d'un pentagone à celles de ses cinq oreilles.

Le cœur du dossier est ailleurs : dans les polygones où l'application
« polygone $\mapsto$ aires des oreilles » perd de l'information même
infinitésimalement, c'est-à-dire où l'on peut déformer le polygone sans que les
aires bougent au premier ordre dans certaines directions. L'outil est celui des
multiplicateurs de Lagrange, et Grothendieck l'écrit lui-même : on cherche les
combinaisons des aires dont la variation est nulle pour toute déformation qui
garde le polygone fermé. Ces combinaisons forment un espace vectoriel, et la
condition s'écrit comme une relation entre quatre côtés consécutifs, une par
sommet. Le théorème de la dernière page classe les cas : en général cet espace
est nul ; il peut être une droite, engendrée par une combinaison où tous les
sommets comptent ; il est un plan exactement quand le polygone est un triangle
parcouru plusieurs fois ; et il peut être une droite engendrée par une
combinaison qui ignore un sommet, ce qui arrive pour des polygones en étoile
d'une forme très particulière, dont une page dessine un exemple à neuf sommets.
Tout cela est conduit en essayant les cas un à un --- triangle, quadrilatère,
pentagone, ennéagone --- avant l'énoncé général, qui arrive à la fin.

Trois pages, écrites à rebours au milieu du dossier, n'ont rien à voir avec les
polygones : elles formalisent une suite de questions posées à une configuration
inconnue, chacune choisie selon les réponses aux précédentes. La première page,
enfin, définit les polygones dans le langage des schémas, comme un espace dont
les points sont des polygones, dégénérés compris.

Les noms à chercher ensuite : multiplicateurs de Lagrange et points critiques,
espaces de modules de polygones, géométrie affine équi-aire, formule de Gauss
pour le pentagone, carré symétrique d'un plan, arbres de décision.

\keywords{plane polygon, signed area, affine geometry, moduli space of polygons, projective frame, Lagrange multipliers, critical points, symmetric square, Gauss pentagon formula, decision tree}
\end{resume}

\pagerange{1}{28}
\section*{Le fil du dossier, et les conventions}

Le dossier n'a pas de titre de sa main ; celui de l'inventaire, « Polygones »,
lui convient. La « lettre » que l'inventaire mentionne n'est sur aucun des
feuillets numérisés, et les pages 2, 7, 11 et 13 sont blanches. Une
pagination de sa main, cerclée, court de 1 à 16 sur les pages 14 à 28 (le 7
manque) : c'est une rédaction suivie, précédée de feuillets d'essais.

\subsection*{Les stations}

\begin{itemize}
  \item \emph{Page 1.} Le foncteur $\mathrm{Polg}_n(E)$ des $n$-gones d'un plan
    affine, sommets et droites des côtés ; dégénérescences.
  \item \emph{Pages 3 et 4.} La relation des multiplicateurs sous sa première
    forme, $\lambda_i u_{i-1} - \lambda_{i+1} u_{i+1} = w$ constant, et ses
    différences $M_i$ ; le calcul des côtés de proche en proche.
  \item \emph{Pages 5 et 6.} Le quadrilatère : une relation linéaire entre les
    aires des oreilles, et l'espace des quadrilatères identifié à un plan
    projectif.
  \item \emph{Pages 8 à 10.} Un morceau étranger : configurations, opérations,
    ensembles résiduels, tests, « cribles ».
  \item \emph{Page 12.} Les relations entre multiplicateurs, aires et le
    vecteur $w$.
  \item \emph{Pages 14 et 15.} La formulation lagrangienne écrite en
    coordonnées, et la formule close de $\lambda_i$.
  \item \emph{Pages 16 et 17.} Le quadrilatère à nouveau : son multiplicateur,
    une carte explicite, les diagonales.
  \item \emph{Page 18.} Un multiplicateur nul en un sommet : l'ennéagone étoilé.
  \item \emph{Page 19.} Le pentagone, et une équation du second degré.
  \item \emph{Pages 20 à 22.} La relation à quatre termes ; $u_i$ dans la base
    $u_0, u_1, u_2$ ; les polynômes $P_i, Q_i, R_i$ ; les équations de
    fermeture.
  \item \emph{Pages 23 à 27.} Le cas dégénéré, la Proposition et sa preuve.
  \item \emph{Page 28.} Le Théorème : les quatre cas.
\end{itemize}

\subsection*{Conventions}

$k$ est un corps, $V$ un plan vectoriel sur $k$, $E$ un plan affine de
direction $V$. Un $n$-gone a pour sommets $s_i$ et pour côtés
$u_i = s_{i+1} - s_i$, $i \in \mathbf{Z}/n\mathbf{Z}$ ; il est \emph{fermé}
quand $\sum u_i = 0$, ce qui est automatique si on le donne par ses sommets,
et ne l'est plus si on le donne par la suite $u_* = (u_i)$ de ses côtés. On
note, comme la page 15,
\[
  u_{ab} = u_a + u_{a+1} + \cdots + u_{b-1} = s_b - s_a .
\]
On fixe une forme volume sur $V$, et l'on pose
\[
  \mu_i = u_{i-1} \wedge u_i \in k ,
\]
le double de l'aire algébrique de l'\emph{oreille} en $s_i$, le triangle
$s_{i-1} s_i s_{i+1}$\footnote{Les pages 5 et 12 écrivent $\mu_i$, les pages 14
à 19 $f_i$ ; on garde $\mu_i$ partout. Le nom d'« oreille » est le nôtre ; les
pages ne nomment pas ces triangles. Une transformation affine de partie linéaire
$g$ multiplie tous les $\mu_i$ par $\det g$ : la classe de $(\mu_i)$ dans
$\mathbf{P}^{n-1}$ est un invariant affine.}. On note $\varphi$ l'application
$u_* \mapsto (\mu_i)_i$.

La relation des multiplicateurs a deux formes, et le dossier passe de l'une à
l'autre. La \emph{forme lagrangienne} (pages 3, 4, 12, 14 et 15) :
\[
  (L) \qquad \lambda_i u_{i-1} - \lambda_{i+1} u_{i+1} = w
  \quad \text{indépendant de } i ;
\]
la \emph{forme à quatre termes} (pages 20 à 28), qui en est la différence de
deux lignes consécutives :
\[
  (\ast) \qquad \lambda_i u_{i-1} - \lambda_{i+1}(u_i + u_{i+1})
  + \lambda_{i+2} u_{i+2} = 0 \qquad \forall\, i \in \mathbf{Z}/n\mathbf{Z} .
\]
Les deux sont équivalentes : $(\ast)$ dit exactement que le membre de gauche de
$(L)$ ne dépend pas de $i$. L'espace des $(\lambda_i) \in
k^{\mathbf{Z}/n\mathbf{Z}}$ qui les satisfont est noté $N(u_*)$, comme à la
page 25\footnote{La page 20 écrit la même relation décalée d'un cran,
$\lambda_{i-1} u_{i-2} - \lambda_i (u_{i-1} + u_i) + \lambda_{i+1} u_{i+1} = 0$ ;
on suit l'indexation des pages 25 à 28. Le vecteur constant $w$ s'appelle $u$
aux pages 3, 4 et 12 et a pour coordonnées $(\nu, -\mu)$ à la page 14 ; on le
renomme pour garder $u$ libre.}.

\subsection*{Ce que seule une lecture d'ensemble peut dire}

Trois raccords, qu'aucune page ne fait.

Le $N(u_*)$ du théorème final est l'espace des multiplicateurs de Lagrange de
la page 14 : la relation à quatre termes des pages 20 à 28 est la différence
$M_i = L_{i-1} - L_i$ que la page 3 forme déjà. Ce qui est classé à la fin est
donc l'ensemble des points où $\varphi$, restreinte aux polygones fermés, n'est
pas une submersion, et de combien elle ne l'est pas.

Le quadrilatère est toujours dans ce cas : la relation entre ses aires
(page 5) a pour coefficients $(1, -1, 1, -1)$, et c'est exactement le
multiplicateur $\lambda_0 + \lambda_1 = 0$ que les pages 16 et 17 retrouvent
par le calcul.

Le critère du cas spécial II, au bas de la page 28, est la formule close des
pages 15 et 18 avec $\lambda_0 = 0$, et l'ennéagone étoilé de la page 18 en
est un exemple : son espace de multiplicateurs est une droite, engendrée par
un vecteur nul en $0$ et en $0$ seulement.

\subsection*{Ce que le dossier annonce sans l'établir}

La condition qui définit l'ouvert $\mathcal{U}$ de la page 1 est biffée ; les
comptes de dimensions de la page 4 ne se lisent plus ; le système des équations
de fermeture (page 22) n'est jamais résolu ; le sens de $CI_n^*$ (page 12)
n'est pas fixé ; la quadrique du haut de la page 20 n'est rattachée à rien ;
la démonstration de d) $\Rightarrow$ a) (page 27) est annoncée et la page est
blanche ensuite ; la note de marge de la page 28 est en grande partie perdue.
Les « cribles » des pages 8 à 10 ne servent nulle part dans le dossier.

\pagerange{1}{1}
\section*{Le foncteur des polygones (page 1)}

Soit $S$ un schéma, $E$ un plan affine sur $S$, et $n \geqslant 3$. Pour un
$S$-schéma $S'$, on pose
\[
  \mathrm{Polg}_n(E)(S') = \bigl\{ (s_0, \ldots, s_{n-1}, d_0, \ldots,
  d_{n-1}) \bigm| s_i \in E(S'),\ d_i \in \mathrm{Dr}(E_{S'}),\
  s_i \in d_i,\ s_{i+1} \in d_i \bigr\},
\]
où $\mathrm{Dr}(E_{S'})$ est l'ensemble des droites de $E_{S'}$ et les indices
sont pris modulo $n$\footnote{Le nom « Polg » est lu avec doute ; son
initiale est doublée comme une capitale ajourée. La page écrit l'incidence par
un signe $\prec$ pour les premiers couples et par $\in$ pour le dernier.}. Un
point est un $n$-gone muni des droites de ses côtés. Les conditions d'incidence
sont fermées, si bien que $\mathrm{Polg}_n(E)$ est représentable par un
sous-schéma fermé de $E^n \times \mathrm{Dr}(E)^n$, le schéma des droites d'un
plan affine étant un ouvert du plan projectif dual\footnote{La représentabilité
est notre remarque ; la page ne dit rien de ce genre.}.

Garder les droites $d_i$ à côté des sommets, c'est ce qui permet de suivre les
dégénérescences. Les croquis de la page en montrent trois : des côtés
consécutifs portés par une même droite, $d_0 = d_1 = \cdots = d_i \neq d_{i+1}$,
les sommets $s_0, \ldots, s_{i+1}$ étant alors alignés ; deux sommets
consécutifs confondus, $s_{i+1} = s_{i+2}$, où le côté nul garde la droite qui
le porte ; un triangle vu comme un hexagone dégénéré, de sommets
$s_0 = s_1$, $s_2 = s_3$, $s_4 = s_5$ et de côtés $d_0 = d_1$, $d_2 = d_3$,
$d_4 = d_5$.

La page définit encore un ouvert $\mathcal{U}$, indexé par un ensemble
$\underline{S}$ de sommets et un ensemble $\underline{A}$ d'arêtes --- un
polygone combinatoire $(S, A)$ au sens que le dossier 89 (pages 19 et 20)
précise ---, dont la condition est biffée\footnote{On lit encore sous la
rature « $P_s \notin \ldots d_{s-2}$ », peut-être un sommet hors de la droite
d'un côté non adjacent ; on ne reconstitue pas la condition. L'indice de
$\mathcal{U}$ est lu avec doute.}.

\pagerange{3}{4}
\section*{La relation des multiplicateurs, première forme (pages 3 et 4)}

Ces deux pages sont des essais. On y lit la relation sous la forme $(L)$ :
pour un $n$-gone donné par ses côtés, des scalaires $\lambda_i$ tels que
\[
  L_i := \lambda_i u_{i-1} - \lambda_{i+1} u_{i+1}
\]
ne dépende pas de $i$ : $L_0 = L_1 = \cdots = L_{n-1} = w$. Les différences
\[
  M_i = L_{i-1} - L_i = \lambda_{i-1} u_{i-2} - \lambda_i (u_{i-1} + u_i)
  + \lambda_{i+1} u_{i+1}
\]
sont les premiers membres de la relation à quatre termes $(\ast)$, qui
reviendra seule aux pages 20 à 28\footnote{La page écrit
$M_{n-1} = L_{n-2} - L_{n-1}$ sous une accolade ; les deux lignes qui suivent,
qui développent $M_1 u_1$ et $M_2 u_2$, s'interrompent sur un « mod » en
suspens. Le produit $u_{-1} u_{12}$ qui y figure est celui de l'algèbre
symétrique, voir la section des pages 14 et 15.}.

Sous le titre « Calcul intrinsèque »\footnote{Lecture incertaine.}, la
page 4 calcule les côtés de proche en proche. Si les $\lambda_i$ sont non nuls,
$(L)$ en $i$ donne $u_{i+1} = (\lambda_i u_{i-1} - w)/\lambda_{i+1}$ : les
données $u_0, u_1, w$ et $\lambda_*$ déterminent tous les $u_i$, et les deux
dernières lignes, $u_n = u_0$ et $u_{n+1} = u_1$, sont deux relations entre
$u_0, u_1, w, \lambda_*$. C'est le schéma que les pages 20 à 22 reprendront
avec la forme à quatre termes. La page écrit aussi la ligne
$\lambda_2 \, u_1 u_2 = -\lambda_0 \, u_{-1} u_1 + \lambda_1 \, u_1 (u_0 + u_1)$,
qui est le cas $i = 2$ de la formule close de la page 15\footnote{La prose du
milieu de la page, qui comptait des dimensions --- « $n-2$
($\subset \mathbf{P}^{n-1}$) », « dans le cas $\sum u_i = 0$ », « codim.~3 »
---, est presque entièrement illisible, de même qu'une note de marge sur des
$(c_1, \ldots, c_{n-2})$ de sommes alternées nulles. Pour ce que valent ces
comptes : les $n$-gones fermés non alignés à transformation affine près forment
un espace de dimension $2n - 6$, et la classe de $(\mu_i)$ vit dans
$\mathbf{P}^{n-1}$ ; les deux dimensions coïncident pour $n = 5$, ce qui est la
situation de la page 19. Ce compte est le nôtre.}.

\pagerange{5}{6}
\section*{Les quadrilatères forment un plan projectif (pages 5 et 6)}

Soit un quadrilatère fermé, $u_0 + u_1 + u_2 + u_3 = 0$, et
$\mu_0 = u_3 \wedge u_0$, $\mu_1 = u_0 \wedge u_1$, $\mu_2 = u_1 \wedge u_2$,
$\mu_3 = u_2 \wedge u_3$. En remplaçant $u_3$ par $-(u_0 + u_1 + u_2)$ on
trouve
\[
  u_0 \wedge u_2 = \mu_0 - \mu_1 = \mu_3 - \mu_2 ,
  \qquad
  u_1 \wedge u_3 = \mu_0 - \mu_2 ,
\]
d'où la relation
\[
  \mu_0 + \mu_2 = \mu_1 + \mu_3 .
\]
Les deux membres sont le double de l'aire du quadrilatère, coupé par l'une ou
l'autre diagonale ; la page 17 vérifie qu'ils valent aussi
$(u_0 + u_1) \wedge (u_1 + u_2)$, le produit extérieur des deux
diagonales\footnote{Page 5, la deuxième ligne du calcul de $u_0 \wedge u_2$
doit se lire $-u_0 \wedge (u_1 + u_3)$ ; le signe est surchargé. Page 17, le
point d'interrogation de « $\overset{?}{=}$ » est levé par le calcul qui
suit, que la page conclut d'un « ok ».}.

Soit $\mathbf{M}_4$ l'ensemble des quadrilatères fermés non contenus dans une
droite, à transformation affine près. L'application
\[
  \mathbf{M}_4 \longrightarrow \bigl\{ \mu_0 + \mu_2 = \mu_1 + \mu_3 \bigr\}
  \subset \mathbf{P}^3, \qquad
  (u_0, u_1, u_2, u_3) \longmapsto (\mu_0 : \mu_1 : \mu_2 : \mu_3)
\]
est bien définie (un quadrilatère non aligné a un $\mu_i$ non nul) et c'est une
bijection sur ce plan projectif. Sur l'ouvert $\mu_1 \neq 0$, en effet,
$u_0, u_1$ forment une base de $V$, qu'une transformation affine envoie sur une
base fixée ; écrivant $u_2 = \alpha u_0 + \beta u_1$, on trouve
$\mu_1 = 1$, $\mu_2 = -\alpha$, $\mu_3 = \beta - \alpha$, $\mu_0 = 1 + \beta$
dans la base $e = u_0 \wedge u_1$, d'où
\[
  u_2 = -\frac{\mu_2}{\mu_1}\, u_0 + \frac{\mu_3 - \mu_2}{\mu_1}\, u_1 ,
  \qquad
  u_3 = \Bigl( \frac{\mu_2}{\mu_1} - 1 \Bigr) u_0 - \frac{\mu_0}{\mu_1}\, u_1 ,
\]
et tout point du plan avec $\mu_1 \neq 0$ est atteint une fois. Les autres
cartes s'en déduisent par permutation circulaire\footnote{La page écrit
« $\mathbf{M}_4 \to \mathbf{P}^2_{\mathbf{Z}}$ mono », « injectif » étant
biffé, puis « épi ? », et conclut « $\mathbf{M}_4 \simeq \mathbf{P}^2_{\mathbf{Z}}$ »
après les formules ci-dessus : la question est réglée par ces formules
elles-mêmes, ce qu'on dit ici. L'indice $\mathbf{Z}$ indique qu'il pense
l'énoncé sur $\mathrm{Spec}\,\mathbf{Z}$ ; on ne l'a vérifié que sur un corps.
Il ne précise pas ce qu'il exclut ; les quadrilatères alignés ont tous leurs
$\mu_i$ nuls et n'ont pas d'image.}.

\emph{Les quadrilatères dégénérés.} Un côté est nul si et seulement si les deux
aires qui le contiennent le sont : $u_i = 0$ équivaut à
$\mu_i = \mu_{i+1} = 0$ (la page encadre le cas $u_2 = 0 \iff \mu_2 = \mu_3 =
0$). Les quatre quadrilatères dégénérés sont donc les points
\[
  (0,0,1,1), \quad (1,0,0,1), \quad (1,1,0,0), \quad (0,1,1,0)
\]
du plan, pour $u_0, u_1, u_2, u_3$ nuls respectivement. Les trois premiers
sont indépendants (la matrice qu'ils forment est de rang 3) et le quatrième est
leur somme alternée, à coefficients tous non nuls : ces quatre points sont en
position générale, ils forment un \emph{repère projectif} du plan. Par
conséquent, si $\Pi = (S, A)$ est un quadrilatère combinatoire --- quatre
sommets, quatre arêtes, au sens du dossier 89 (pages 19 et 20) ---, l'espace
$\mathbf{M}(\Pi)$ des quadrilatères de type $\Pi$ s'identifie canoniquement au
plan projectif engendré par $A$, c'est-à-dire au plan
$\mathbf{P}\bigl(k^{A} / k\cdot(1,1,1,1)\bigr)$ dont les quatre points de
repère sont les arêtes, chaque arête correspondant au quadrilatère où elle est
réduite à un point\footnote{La page dit « l'enveloppe projective de $A$ »,
« enveloppe » étant lu avec doute ; la formule
$\mathbf{P}(k^{A}/k\cdot(1,1,1,1))$ est la nôtre. La page 6 ne nomme que trois
des quatre points ; leur « rang 3 » et le mot « coplanaires » sont des lectures
douteuses, la phrase étant en partie illisible. Le dossier 75 (pages 17 à 19)
associe de même à un polygone combinatoire un plan canonique, euclidien
celui-là, où il est réalisé régulièrement ; les deux constructions ne se
rejoignent pas.}. Cinq croquis de quadrilatères --- ouvert, croisé, aplati,
avec ou sans leurs diagonales --- terminent la page.

\pagerange{8}{10}
\section*{Un morceau étranger : opérations, tests et cribles (pages 8 à 10)}

Ces trois pages ne parlent pas de polygones, et ne disent pas à quoi elles
servent. Elles sont écrites à rebours : la page 10 pose les définitions, la
page 9 construit l'objet, la page 8 le récrit. On les lit dans l'ordre logique.

\emph{Les données} (page 10). Un ensemble $C$ de « configurations », un
ensemble $\Omega$ d'« opérations » (virtuelles), et pour chaque
$\omega \in \Omega$ un ensemble $V_\omega$ de « résultats » possibles et une
application $f_\omega : C \to V_\omega$, notée $x \mapsto \omega(x)$ : le
résultat de l'opération $\omega$ sur la configuration $x$. Une
\emph{opération effectuée} est une opération avec son résultat,
\[
  \widetilde{\Omega} = \coprod_{\omega \in \Omega} V_\omega
  \xrightarrow{\ \pi\ } \Omega ,
  \qquad \widetilde{\omega} = (\omega, v) .
\]
Son \emph{ensemble résiduel} est $C(\widetilde{\omega}) = f_\omega^{-1}(v)$,
l'ensemble des configurations compatibles avec ce résultat, et pour une suite,
$C(\widetilde{\omega}_1, \ldots, \widetilde{\omega}_\nu) = \bigcap_i
C(\widetilde{\omega}_i)$. Chaque configuration $x$ est une section de $\pi$,
$\omega \mapsto (\omega, \omega(x))$, d'où des applications
$C \times \Omega^\nu \to \widetilde{\Omega}^\nu$. Les \emph{tests effectués de
longueur $\nu$} forment l'image
\[
  T_\nu = \mathrm{Im}\bigl( C \times \Omega^\nu \to \widetilde{\Omega}^\nu
  \bigr) \subset \widetilde{\Omega}^\nu ,
\]
c'est-à-dire les suites de résultats que produit effectivement une même
configuration --- autrement dit, celles dont l'ensemble résiduel n'est pas
vide\footnote{Cette dernière reformulation est la nôtre. Les mots
« opérations », « résiduel » et « effectués » sont lus avec doute ; un NB de la
page 10 est illisible.}.

\emph{Les cribles} (page 9). Un \emph{crible de longueur $\nu$} est une suite
de choix : une partie non vide $\mathcal{F}_1 \subset \Omega$ des premières
opérations permises ; puis, pour chaque première opération effectuée possible,
$\widetilde{\omega}_1 \in \widetilde{\mathcal{F}}_1 = \pi^{-1}(\mathcal{F}_1)
\cap T_1$, une partie $f_2(\widetilde{\omega}_1) \subset \Omega$ des secondes
opérations permises ; et ainsi de suite,
\[
  \widetilde{\mathcal{F}}_i = \bigl\{ (\widetilde{\omega}_1, \ldots,
  \widetilde{\omega}_i) \in T_i \bigm| (\widetilde{\omega}_1, \ldots,
  \widetilde{\omega}_{i-1}) \in \widetilde{\mathcal{F}}_{i-1},\
  \omega_i \in f_i(\widetilde{\omega}_1, \ldots, \widetilde{\omega}_{i-1})
  \bigr\},
  \qquad
  f_i : \widetilde{\mathcal{F}}_{i-1} \to \mathfrak{P}(\Omega),
\]
jusqu'à $\widetilde{\mathcal{F}}_\nu \subset T_\nu$. Certains
$\widetilde{\mathcal{F}}_i$ peuvent être vides\footnote{« Crible de longueur
$\nu$ » est écrit au-dessus d'un mot biffé, et lu avec doute. La page encadre
$f_2 : \widetilde{\mathcal{F}}_1 \to \mathfrak{P}^*(\Omega)$, l'astérisque
désignant sans doute les parties non vides, puis écrit $\mathfrak{P}(\Omega)$
pour les $f_i$ suivants. Elle écrit
« $\widetilde{\omega}_2 \in f_2(\widetilde{\omega}_1)$ » là où c'est
l'opération sous-jacente $\omega_2$ qui doit être permise.}.

\emph{Les graphes} (page 8). La même donnée se récrit en remplaçant chaque
$f_i$ par son graphe, une partie de $T_{i-1} \times \Omega \subset
\widetilde{\Omega}^{i-1} \times \Omega$, et l'on a
$\mathcal{F}_i \subset T_i \subset \widetilde{\Omega}^i$, la projection de
$\mathcal{F}_i$ sur les $i-1$ premiers facteurs tombant dans
$\mathcal{F}_{i-1}$.

C'est ce qu'on appelle aujourd'hui une stratégie adaptative de questionnement,
ou un arbre de décision : on interroge une configuration inconnue, chaque
question étant choisie selon les réponses précédentes, et les ensembles
résiduels mesurent ce qui reste à savoir. Le crible est ici non déterministe
--- un ensemble de questions permises à chaque étape plutôt qu'une seule ---, et
rien n'est encore dit de son but\footnote{Les noms sont les nôtres.}.

\pagerange{12}{12}
\section*{Relations entre multiplicateurs et aires (page 12)}

Sous les hypothèses
\[
  \lambda_i u_{i-1} - \lambda_{i+1} u_{i+1} = w \quad (i \in \mathbf{Z}/n\mathbf{Z}),
  \qquad \sum u_i = 0 ,
\]
posons $\mu_i = u_{i-1} \wedge u_i$ et $\alpha_i = w \wedge u_i$. En faisant
le produit extérieur de $(L)$ avec $u_i$, avec $w$, et en sommant :
\[
  \lambda_i \mu_i + \lambda_{i+1} \mu_{i+1} = \alpha_i ,
  \qquad
  \lambda_i \alpha_{i-1} - \lambda_{i+1} \alpha_{i+1} = 0 ,
  \qquad
  \sum \alpha_i = 0 ,
\]
et donc, sommant la première sur $i$,
\[
  \sum_i \lambda_i \mu_i = 0 .
\]
Cette dernière identité dit qu'un multiplicateur s'annule sur le vecteur
$\varphi(u_*)$ lui-même : c'est la trace de l'homogénéité de $\varphi$, qui est
de degré 2\footnote{L'interprétation est la nôtre ; la page écrit
« $2 \sum \lambda_i \mu_i = 0$ ».}.

La page note que $u_0, u_1, w$ et $(\lambda_i)$ déterminent tous les $u_i$
(c'est le calcul de la page 4), puis que $\alpha_0, \alpha_1$ et $(\lambda_i)$
les déterminent à transformation affine près si $\mu_1 \neq 0$ : décomposant
$w$ dans la base $u_0, u_1$,
\[
  w = \frac{\alpha_1}{\mu_1}\, u_0 - \frac{\alpha_0}{\mu_1}\, u_1
  = \Bigl( \lambda_1 + \lambda_2 \frac{\mu_2}{\mu_1} \Bigr) u_0
  - \Bigl( \lambda_0 \frac{\mu_0}{\mu_1} + \lambda_1 \Bigr) u_1 ,
\]
et $(L)$ en $i = 1$ donne
\[
  u_2 = -\frac{\mu_2}{\mu_1}\, u_0 + \frac{1}{\lambda_2}
  \Bigl( \lambda_0 \frac{\mu_0}{\mu_1} + \lambda_1 \Bigr) u_1 ,
\]
et ainsi de suite jusqu'à $u_{n-1} = A_{n-1}\, u_0 + B_{n-1}\, u_1$, les
coefficients étant des fonctions de $(\lambda)$ et $(\mu)$. Les couples
$(\lambda_i, \mu_i)_i$ sont considérés à un scalaire près, et la page esquisse
une application
\[
  CI_n^* \longrightarrow \mathbf{P}^{n-1} \times \mathbf{P}^{n-1} ,
  \qquad (\lambda), (\mu) ,
\]
où $CI_n^*$ désigne des « $n$-polygones cycliques (mod Aff) »\footnote{Le
reste de la définition de $CI_n^*$ est illisible. Ce qui précède suggère qu'il
s'agit des $n$-gones munis d'un multiplicateur, envoyés sur le couple formé du
multiplicateur et des aires ; on ne l'affirme pas. La colonne de droite,
séparée par un trait et très raturée, énumère des hypothèses :
$(\mu_i)_i \neq 0$ et $w \neq 0$ si et seulement si $(\alpha_i)_i \neq 0$ ;
$u_i \neq 0$, $\lambda_i \neq 0$ pour tout $i$ ; $(\lambda_i, \alpha_i)_i \neq 0$ ;
et un « $\mathbf{P}(k^n \times L^n)$ », la lettre $L$ étant douteuse.}.

\pagerange{14}{15}
\section*{Les multiplicateurs de Lagrange, et la formule close (pages 14 et 15)}

La page 14, première de la rédaction suivie, dit d'où vient la relation.
Écrivons les côtés en coordonnées, $u_i = (x_i, y_i)$, de sorte que
$\mu_i = x_{i-1} y_i - y_{i-1} x_i$. Alors
\[
  \sum_i \lambda_i \, d\mu_i
  = \sum_i \bigl( -\lambda_i y_{i-1} + \lambda_{i+1} y_{i+1} \bigr)\, dx_i
  + \sum_i \bigl( \lambda_i x_{i-1} - \lambda_{i+1} x_{i+1} \bigr)\, dy_i ,
\]
et cette forme est une combinaison $a \sum dx_i + b \sum dy_i$ des
différentielles de la condition de fermeture si et seulement si les
coefficients de $dx_i$ et de $dy_i$ ne dépendent pas de $i$, c'est-à-dire si et
seulement si $(L)$ a lieu. Autrement dit :

\textbf{La relation $(L)$, ou $(\ast)$, dit que $(\lambda_i)$ est un
multiplicateur de Lagrange.} $N(u_*)$ est l'espace des formes linéaires sur
$k^{\mathbf{Z}/n\mathbf{Z}}$ qui annulent l'image de la différentielle de
$\varphi$ en $u_*$, restreinte aux déformations $(du_i)$ telles que
$\sum du_i = 0$ : l'espace conormal à l'image, en ce point, de l'application
« aires des oreilles » sur l'espace des polygones fermés. En particulier
$N(u_*) = 0$ si et seulement si cette application est une submersion en
$u_*$\footnote{La page écrit l'égalité
$\sum \lambda_i df_i = \mu \sum dx_i + \nu \sum dy_i$ et en tire les deux
systèmes de relations ; les mots « multiplicateur de Lagrange », « conormal »
et « submersion » sont les nôtres. Le calcul ne suppose pas $\sum u_i = 0$ :
pour une suite $u_*$ quelconque, $N(u_*)$ est le conormal de la restriction de
$\varphi$ à l'hyperplan affine $\sum u_i = \mathrm{cste}$ qui passe par $u_*$.}.

\emph{La formule close.} Supposons $\sum u_i = 0$ et toutes les coordonnées
$x_i, y_i$ non nulles\footnote{« OPS $x_i, y_i$ tous $\neq 0$ (par choix de
coordonnées) » : c'est possible dès que $k$ a assez d'éléments, par exemple
s'il est infini. L'hypothèse ne sert qu'à diviser ; l'identité finale ne la
demande pas.}. La première coordonnée de $(L)$ donne
$\lambda_{i+1} y_{i+1} = \lambda_i y_{i-1} - c$ avec $c$ constant, et une
récurrence immédiate sur $i$ donne, pour $2 \leqslant i \leqslant n-1$,
\[
  (y_{i-1} y_i)\, \lambda_i = (-\lambda_0 y_{-1})(y_1 + \cdots + y_{i-1})
  + (\lambda_1 y_1)(y_0 + \cdots + y_{i-1}) ,
\]
et la même formule pour les $x$. Pour $i = n$ elle se réduit à l'identité
$\lambda_0 y_{-1} y_0 = \lambda_0 y_0 y_{-1}$, grâce à
$y_1 + \cdots + y_{n-1} = -y_0$ et $y_0 + \cdots + y_{n-1} = 0$ ; la page
conclut « ok »\footnote{La page 14 commence le calcul avec des produits de
plus en plus longs, $(y_2 y_3 y_4)\lambda_4$, $(y_2 y_3 y_4 y_5)\lambda_5$,
dont les indices sont surchargés ; la forme encadrée de la page 15, à deux
facteurs seulement, est celle que la récurrence établit, et c'est la seule
qu'on garde.}.

Comme la même identité vaut pour toute forme linéaire $\ell$ sur $V$ à la place
de la coordonnée $y$ --- la récurrence est la même ---, elle est une identité
entre formes quadratiques sur $V^*$, c'est-à-dire dans le carré symétrique
$\operatorname{Sym}^2 V$ :
\[
  \lambda_i \; u_{i-1} \cdot u_i
  = -\lambda_0 \; u_{-1} \cdot u_{1i} + \lambda_1 \; u_1 \cdot u_{0i} ,
  \qquad u_{ab} = s_b - s_a ,
\]
où $\cdot$ est le produit de l'algèbre symétrique\footnote{La page écrit cette
forme vectorielle, encadrée, avec une simple juxtaposition
« $\lambda_i u_{i-1} u_i = -\lambda_0 u_{-1} u_{1i} + \lambda_1 u_1 u_{0i}$ ».
L'en-tête de la transcription rapproche cette juxtaposition du déterminant ;
elle ne peut pas l'être : avec $u_{i-1} \wedge u_i$, la relation $(L)$ en
$i = 1$ donne pour $i = 2$ le terme $+\lambda_0\, u_{-1} \wedge u_1$, de signe
opposé à celui de la formule. Ce que la page établit, coordonnée par
coordonnée, est l'identité dans $\operatorname{Sym}^2 V$ qu'on écrit ici ; le
nom est le nôtre.}. Éliminant $\lambda_i$ entre la formule en $x$ et la formule
en $y$, la page en tire, pour $2 \leqslant i \leqslant n-1$, une relation
linéaire entre $\lambda_0$ et $\lambda_1$ seuls :
\begin{align*}
  &\lambda_0 \bigl[ -x_{i-1} x_i y_{-1} (y_1 + \cdots + y_{i-1})
  + y_{i-1} y_i x_{-1} (x_1 + \cdots + x_{i-1}) \bigr] \\
  &\quad + \lambda_1 \bigl[ x_{i-1} x_i y_1 (y_0 + \cdots + y_{i-1})
  - y_{i-1} y_i x_1 (x_0 + \cdots + x_{i-1}) \bigr] = 0 .
\end{align*}
Un essai sur le triangle, très raturé, suit\footnote{On n'en restitue rien. En
marge, pour le triangle $u, v, w = -(u+v)$, la page note que
$u \wedge v = v \wedge w = w \wedge u$ : les trois oreilles d'un triangle sont
le triangle lui-même.}.

\pagerange{16}{17}
\section*{Le quadrilatère, deuxième passage (pages 16 et 17)}

Pour $n = 4$ ($y_{-1} = y_3$, $x_{-1} = x_3$), la relation de la page 15 en
$i = 2$ s'écrit, les deux crochets étant égaux,
\[
  (\lambda_0 + \lambda_1) \cdot \bigl( -x_1 y_1\, (u_2 \wedge u_3) \bigr) = 0 ,
\]
et en $i = 3$
\[
  (\lambda_0 + \lambda_1) \cdot x_3 y_3\, (u_1 \wedge u_2) = 0 .
\]
Donc $\lambda_0 + \lambda_1 = 0$ dès que le sommet $s_3$, ou le sommet $s_2$,
n'est pas aplati\footnote{« Aplati » est lu avec doute. Pour $i = 3$ la page
écrit le crochet en $\lambda_0$ avec le signe opposé, $-x_3 y_3 (x_1 y_2 - y_1
x_2)$ (le signe du second résultat est surchargé), et rattache la condition au
sommet $s_1$, avec un croquis de l'angle entre $u_1$ et $u_2$ ; cet angle est
celui du sommet $s_2$. La conclusion encadrée, $\lambda_0 + \lambda_1 = 0$, est
juste dans les deux cas.}. Par permutation circulaire, si le quadrilatère n'est
pas aligné, $\lambda_i + \lambda_{i+1} = 0$ pour tout $i$ : $N(u_*)$ est la
droite engendrée par $(1, -1, 1, -1)$, qui est bien dans $N(u_*)$ puisque
$(\ast)$ se réduit alors à $\lambda_i \sum_j u_j = 0$. C'est exactement la
relation $\mu_0 - \mu_1 + \mu_2 - \mu_3 = 0$ de la page 5 : l'image de
$\varphi$ est contenue dans cet hyperplan, et son équation est le
multiplicateur.

\emph{Une carte explicite} (page 17). Le quadrilatère de sommets
\[
  s_0 = (a - 1, 0), \quad s_1 = (0, b - 1), \quad
  s_2 = (a + 1, 0), \quad s_3 = (0, b + 1)
\]
a ses deux diagonales sur les axes, chacune de longueur 2. On trouve
\[
  \mu_0 = 2(1 - a), \quad \mu_1 = 2(1 - b), \quad
  \mu_2 = 2(1 + a), \quad \mu_3 = 2(1 + b),
\]
donc $\mu_1 + \mu_3 = \mu_0 + \mu_2 = 4$, $\mu_3 - \mu_1 = 4b$,
$\mu_2 - \mu_0 = 4a$\footnote{La page écrit $v$ par-dessus un « 4 » pour la
valeur commune des deux sommes, et « $(b)\,v$ », « $(a)\,v$ » pour les deux
différences.}. Tout quadrilatère dont les diagonales ne sont pas parallèles se
ramène à celui-là par une transformation affine, et $(a, b)$ se lit sur ses
aires : c'est une carte affine $\mathbf{A}^2$ de $\mathbf{M}_4$, le
complémentaire de la droite $\mu_0 + \mu_2 = 0$ des quadrilatères d'aire
nulle\footnote{Cette interprétation est la nôtre ; la page donne le calcul. Sous
une ligne ondulée, elle dessine un pentagone avec toutes ses diagonales, qui
annonce la page 19.}.

\pagerange{18}{18}
\section*{Un multiplicateur nul en un sommet : l'ennéagone étoilé (page 18)}

Si $\lambda_0 = 0$, et en normalisant $\lambda_1 = 1$, la formule close de la
page 15 devient
\[
  \lambda_i \; u_{i-1} \cdot u_i = u_1 \cdot u_{0i} \qquad
  \text{dans } \operatorname{Sym}^2 V .
\]
Dans un plan, le produit de deux vecteurs non nuls de
$\operatorname{Sym}^2 V$ détermine la paire de droites qui les portent ; la
formule dit donc que la paire de directions $\{u_{i-1}, u_i\}$ est la paire
$\{u_1, s_i - s_0\}$, chaque fois que $s_i \neq s_0$. Le dessin de la page en
donne un exemple, un ennéagone étoilé dont on joint le sommet $s_0$ à tous les
autres, en coordonnées :
\[
  \begin{array}{lll}
  s_0 = (0,0), & s_1 = (1,0), & s_2 = (1,1), \\
  s_3 = (\beta_1, \beta_1), & s_4 = (\beta_1, y_4), &
  s_5 = (\beta_2 \beta_1, \beta_2 y_4), \\
  s_6 = (\beta_2 \beta_1, y_6), & s_7 = (\beta_3 \beta_2 \beta_1, \beta_3 y_6), &
  s_8 = (\beta_3 \beta_2 \beta_1, y_8),
  \end{array}
\]
avec $\beta_j \neq 0, 1$ et $y_4 \neq \beta_1$, $y_6 \neq \beta_2 y_4$,
$y_8 \neq \beta_3 y_6$, et $s_9 = s_0$ (ce qui revient à poser $\beta_4 = 0$).
Les côtés d'indice impair sont parallèles à $u_1$, verticaux ; ceux d'indice
pair sont portés par la droite qui joint $s_0$ à leur origine. On vérifie que
$N(u_*)$ est une droite, engendrée par un multiplicateur nul en $0$ et non nul
ailleurs, pour des valeurs génériques des paramètres\footnote{La vérification est la nôtre, par le critère de la page 28
démontré plus bas, et par le calcul sur des valeurs numériques des
paramètres. La page écrit « Aucun des $\lambda_i$ nuls », suivi d'un
« $\iff$ les $s_i$ tous distincts » qu'elle biffe : la bonne condition, on le
verra page 28, est $s_i \neq s_0$ pour $i \neq 0$.}.

La page commence ensuite le calcul général des $u_i$ par la forme à quatre
termes :
\[
  \lambda_3 u_3 = -\lambda_1 u_0 + \lambda_2 u_1 + \lambda_2 u_2 , \quad
  \lambda_4 u_4 = -\lambda_1 u_0 + (\lambda_2 + \lambda_3) u_2 ,
\]
\[
  \lambda_3 \lambda_5 u_5 = -\lambda_1 (\lambda_3 + \lambda_4)\, u_0
  + \lambda_2 \lambda_4\, u_1 + \lambda_2 (\lambda_3 + \lambda_4)\, u_2 ,
\]
que les pages 20 et 21 généralisent\footnote{La page n'écrit pas la valeur de
$\lambda_4 u_4$ au-delà de son premier terme ; on la complète par la relation
qu'elle vient d'écrire, et elle concorde avec le calcul de $\lambda_3 \lambda_5
u_5$ qui suit sur la page.}.

\pagerange{19}{19}
\section*{Le pentagone, et la formule de Gauss (page 19)}

Pour $n = 5$, l'espace des pentagones à transformation affine près et
$\mathbf{P}^4$ ont même dimension. La page normalise
\[
  s_0 = (0,0), \quad s_1 = (1,0), \quad s_2 = (x, y), \quad
  s_3 = (x', y'), \quad s_4 = (0, \mu_0),
\]
ce qui donne $\mu_1 = y$, $\mu_4 = \mu_0 x'$,
$\mu_2 = (x y' - y x') - (y' - y)$ et $\mu_3 = (x y' - y x') + \mu_0 (x' - x)$.
D'où $y = \mu_1$, $x' = \mu_4/\mu_0$, puis
\[
  y' = \mu_0 x + (\mu_1 - \mu_2 + \mu_3 - \mu_4)
\]
et, en reportant, l'équation du second degré
\[
  \mu_0 x^2 - X x + \Bigl( \mu_4 - \mu_3 - \frac{\mu_1 \mu_4}{\mu_0} \Bigr) = 0 ,
  \qquad X = \mu_0 - \mu_1 + \mu_2 - \mu_3 + \mu_4 ,
\]
de discriminant
\[
  \Delta = X^2 - 4\mu_0\mu_4 + 4\mu_0\mu_3 + 4\mu_1\mu_4
  = \sum_i \mu_i^2 - 2 \sum_i \mu_i \mu_{i+1} + 2 \sum_i \mu_i \mu_{i+2}
\]
(indices modulo 5), qu'on peut aussi écrire
\[
  \Delta = (\mu_0 - \mu_1 + \mu_2 + \mu_3 - \mu_4)^2
  + 4 (\mu_1 \mu_3 - \mu_2 \mu_3 + \mu_2 \mu_4) .
\]
Deux pentagones ont donc en général les mêmes aires d'oreilles ; sur
$\mathbf{R}$, ils sont réels tous deux si $\Delta > 0$\footnote{La page
écrit $X = 1 - X_1^4$, avec $X_1^4 = \mu_1 - \mu_2 + \mu_3 - \mu_4$, comme si
$\mu_0 = 1$, puis $X = \mu_0 - X_1^4$, qu'on garde. Les signes des dix doubles
produits de $\Delta$, encadrés un à un, sont surchargés ; la forme
cycliquement symétrique est celle qu'on écrit en premier, sans doute de
lecture, et elle est juste. La ligne médiane,
« $\sum f_i^2 + 2 \sum f_{i+1}(f_i - f_{i+1})$ », lue avec doute, n'est pas
une identité. Dans la dernière ligne la page a $+4 f_2 f_3$ ; il faut
$-4\mu_2\mu_3$.}.

Le calcul cache une formule classique. Notant $A = \sum_i s_i \wedge s_{i+1}$
le double de l'aire du pentagone, on a, dans cette normalisation,
$\mu_0 x = A - \mu_1 - \mu_3$, et l'équation de la page équivaut à
\[
  A^2 - \Bigl( \sum_i \mu_i \Bigr) A + \sum_i \mu_i \mu_{i+1} = 0 ,
\]
dont $\Delta$ est le discriminant : c'est la relation entre l'aire d'un
pentagone et celles de ses cinq oreilles qu'on attribue à Gauss\footnote{Le
rapprochement est le nôtre ; la page ne cite personne, et cette formule n'y
apparaît pas sous cette forme. L'attribution à Gauss est donnée de mémoire. La
dernière ligne de la page, séparée par un trait,
$u_3 \wedge u_4 - u_2 \wedge u_3 = u_3 \wedge (u_2 + u_4)$, s'interrompt.}.

\pagerange{20}{22}
\section*{La forme à quatre termes, et les équations de fermeture (pages 20 à 22)}

La page 20 s'ouvre sur un calcul isolé,
\[
  x^2 + y^2 + z^2 + 2yz + 2zx - 2xy
  = \xi^2 + \eta^2 - 2\xi\eta + 2(\xi + \eta) + 1
  = (\xi - \eta)^2 + 2(\xi + \eta) + 1 ,
\]
où l'on a fait $z = 1$, $x = \xi$, $y = \eta$ : une conique qui, dans la carte
$z = 1$, est la parabole $U^2 + 2V + 1 = 0$, $U = \xi - \eta$,
$V = \xi + \eta$\footnote{La page ne dit pas d'où vient cette forme
quadratique, ni à quoi elle sert ; on ne le devine pas.}. Puis elle encadre la
relation à quatre termes et prend $u_0, u_1$ pour base, $u_2 = \xi u_0 + \eta
u_1$.

\emph{Les côtés dans la base $u_0, u_1, u_2$.} Pour des $\lambda_i$ non nuls,
$(\ast)$ permet de calculer $u_{i+1}$ à partir de $u_{i-2}, u_{i-1}, u_i$ :
\[
  \lambda_{i+1} u_{i+1} = \lambda_i (u_{i-1} + u_i) - \lambda_{i-1} u_{i-2} .
\]
On écrit
\[
  u_3 = \lambda_3^{-1} (P_3 u_0 + Q_3 u_1 + R_3 u_2), \qquad
  P_3 = -\lambda_1,\ Q_3 = R_3 = \lambda_2 ,
\]
\[
  u_i = \frac{\lambda_{i-1}}{\lambda_3 \lambda_4 \cdots \lambda_i}
  \bigl( P_i u_0 + Q_i u_1 + R_i u_2 \bigr) \qquad (i \geqslant 4),
\]
où $P_i, Q_i, R_i \in \mathbf{Z}[\lambda_1, \ldots, \lambda_{i-1}]$ sont des
formes de degré $i - 3$ ; ainsi $(P_4, Q_4, R_4) = (-\lambda_1, 0, \lambda_2 +
\lambda_3)$ et $(P_5, Q_5, R_5) = (-\lambda_1(\lambda_3 + \lambda_4),
\lambda_2\lambda_4, \lambda_2(\lambda_3 + \lambda_4))$, comme à la page 18. En
reportant dans la relation, elles satisfont, chacune, la récurrence à trois
termes
\[
  P_{i+1} = \lambda_{i-1} P_i + \lambda_i \lambda_{i-2} P_{i-1}
  - \lambda_{i-1}^2 \lambda_{i-3} P_{i-2} \qquad (i \geqslant 6),
\]
et de même pour $Q$ et $R$\footnote{Les trois récurrences sont encadrées. La
page écrit d'abord « $5 \leqslant i$ » et le corrige en $i \geqslant 6$ : c'est
juste, la récurrence demande que $u_{i-2}$ soit déjà de la forme générale, donc
$i - 2 \geqslant 4$ ; en $i = 5$ elle est fausse. L'exposant $-1$ de
$\lambda_3$ est lu avec doute ; la page 20 écrivait une première forme, avec
$1/(\lambda_3 \cdots \lambda_i)$ en facteur, que la page 21 remplace.}.

\emph{La fermeture.} Si de plus $\sum u_i = 0$, multiplier par
$\lambda_3 \cdots \lambda_{n-1}$ donne
\[
  (0) \qquad \widetilde{P}_{n-1}\, u_0 + \widetilde{Q}_{n-1}\, u_1
  + \widetilde{R}_{n-1}\, u_2 = 0 ,
\]
où
\[
  \widetilde{P}_{n-1} = \lambda_3 \cdots \lambda_{n-1}
  + \lambda_4 \cdots \lambda_{n-1}\, P_3
  + \sum_{i=4}^{n-1} \lambda_{i-1} \Bigl( \prod_{j=i+1}^{n-1} \lambda_j \Bigr) P_i ,
\]
et de même pour $\widetilde{Q}, \widetilde{R}$, polynômes à coefficients
entiers en $\lambda_1, \ldots, \lambda_{n-1}$\footnote{La page écrit la somme
terme à terme, « $\lambda_4 \lambda_5 \cdots \lambda_{n-1} P_3 + \lambda_3
\lambda_5 \lambda_6 \cdots \lambda_{n-1} P_4 + \lambda_4 \lambda_6 \lambda_7
\cdots \lambda_{n-1} P_5 + \cdots$ », après avoir biffé un début de somme
$\sum_{i=3}^{n-1} \lambda_{i-1} \lambda_{i+1} \cdots$, qui donnerait au terme
$i = 3$ un facteur $\lambda_2$ de trop ; le $\lambda_4$ du troisième terme,
écrit sur un $\lambda_5$, est le bon. Elle écrit d'abord les arguments
$(\lambda_1, \ldots, \lambda_{n-2})$, puis $(\lambda_1, \ldots,
\lambda_{n-1})$ ; c'est la seconde liste qui est juste, $\lambda_{n-1}$
figurant dans les produits. La relation (0), encadrée, porte un numéro lu avec
doute.}. Les trois conditions de périodicité $u_n = u_0$, $u_{n+1} = u_1$,
$u_{n+2} = u_2$, les indices des $\lambda$ étant pris modulo $n$
($\lambda_n = \lambda_0$, $\lambda_{n+1} = \lambda_1$, $\lambda_{n+2} =
\lambda_2$), donnent
\[
  \begin{array}{ll}
  (1) & \bigl[ \lambda_{n-1} P_n - \lambda_3 \cdots \lambda_n \bigr] u_0
    + \lambda_{n-1} Q_n\, u_1 + \lambda_{n-1} R_n\, u_2 = 0 , \\[2pt]
  (2) & \lambda_n P_{n+1}\, u_0
    + \bigl[ \lambda_n Q_{n+1} - \lambda_3 \cdots \lambda_{n+1} \bigr] u_1
    + \lambda_n R_{n+1}\, u_2 = 0 , \\[2pt]
  (3) & \lambda_{n+1} P_{n+2}\, u_0 + \lambda_{n+1} Q_{n+2}\, u_1
    + \bigl[ \lambda_{n+1} R_{n+2} - \lambda_3 \cdots \lambda_{n+2} \bigr] u_2 = 0 ,
  \end{array}
\]
soit un système
\[
  A_j^n\, u_0 + B_j^n\, u_1 + C_j^n\, u_2 = 0 \qquad (j = 0, 1, 2, 3),
\]
avec $A_0^n, B_0^n, C_0^n \in \mathbf{Z}[\Lambda_1, \ldots, \Lambda_{n-1}]
\subset \mathbf{Z}[\Lambda_0, \ldots, \Lambda_{n-1}] \ni A_j^n, B_j^n, C_j^n$
pour $j = 1, 2, 3$, les $\Lambda_i$ étant des indéterminées\footnote{Les trois
relations sont encadrées. Dans (3) la page n'écrit entre $R_{n+2}$ et le
produit qu'un « $(-)$ » ; on rétablit la forme que (1) et (2) imposent. Les
exposants $n$ ne sont nets qu'à la première ligne du système. Page 23, la
page écrit l'inclusion verticalement, par un $\cup$, et donne pour $A_0^n$ les
variables $\Lambda_1, \ldots, \Lambda_{n-2}$ ; voir la note précédente.}. Ce
système, qui exprime qu'une suite de multiplicateurs non nuls existe pour un
polygone fermé donné par $u_0, u_1, u_2$, n'est pas résolu : la page 23 prend
une autre voie.

\pagerange{23}{25}
\section*{Le cas dégénéré, et la Proposition (pages 23 à 25)}

\emph{Le nombre de multiplicateurs.} La relation $(\ast)$ en $i$ fait intervenir
$\lambda_i, \lambda_{i+1}, \lambda_{i+2}$ et donne $\lambda_{i+2} u_{i+2}$ en
fonction des deux autres. Si tous les $u_i$ sont non nuls, deux composantes
consécutives d'un multiplicateur déterminent donc toutes les autres, et
\[
  N(u_*) \longrightarrow k^2, \qquad (\lambda_i) \longmapsto (\lambda_1, \lambda_2)
\]
est injective : $\dim N(u_*) \leqslant 2$\footnote{La page fait le compte avec
trois composantes : « si $\lambda_1, \lambda_2, \lambda_3$ fixés \ldots\ et si
de plus on impose les $\lambda_i \neq 0$, alors les $\lambda_i$
($4 \leqslant i \leqslant n-1$) sont déterminés de façon unique ». C'est vrai,
et le reste pour deux, sans hypothèse sur les $\lambda_i$ : il suffit que les
$u_i$ soient non nuls. On garde l'énoncé le plus fort, dont la démonstration
des pages 26 et 27 se simplifie.}. La relation en $i = 1$,
\[
  \lambda_1 u_0 - \lambda_2 (u_1 + u_2) + \lambda_3 u_3 = 0 ,
\]
montre de plus que si $u_0$, $u_1 + u_2$, $u_3$ ne sont pas colinéaires, les
$(\lambda_1, \lambda_2, \lambda_3)$ possibles forment un espace de dimension 1,
donc $N(u_*)$ aussi, au plus. Par permutation circulaire il en est de même dès
qu'il existe $i$ tel que $u_i$, $u_{i+1} + u_{i+2}$, $u_{i+3}$ ne soient pas
colinéaires. « Mais quand [ils sont] coll.\ $\forall i$ ! »\footnote{La phrase
finale est en partie illisible ; « à la multiplication près » est lu avec
doute. Les crochets de la citation sont de l'édition.}

\emph{Le cas où tous ces triplets sont colinéaires.} Supposons $u_0, u_1$
linéairement indépendants\footnote{« OPS $u_0, u_1$ lin.\ indép. », en marge de
la page 23. On le suppose pour tout le reste de la section, sous la forme
équivalente : les $u_i$ ne sont pas tous sur une même droite (deux côtés
consécutifs sont alors indépendants, et l'on renumérote). La Proposition qui
suit l'omet, et elle est fausse sans lui : voir plus bas.}, et que pour tout
$i$, $u_i$, $u_{i+1} + u_{i+2}$ et $u_{i+3}$ soient colinéaires. Pour $i = 0$ :
$u_1 + u_2 = -\alpha_0 u_0$ et $u_3 = \beta_0 u_0$ ; pour $i = 1$ :
$u_2 + u_3 = \gamma u_1$. Éliminant $u_2$ et $u_3$ dans la base $u_0, u_1$, on
trouve $\gamma = -1$ et $\beta_0 = \alpha_0$, soit
\[
  u_2 = -\alpha_0 u_0 - u_1, \qquad u_3 = \alpha_0 u_0 ,
\]
avec $\alpha_0 \neq 0$ puisque $u_3 \neq 0$. En particulier $u_1, u_2$ sont
indépendants, et de proche en proche tous les couples $(u_i, u_{i+1})$ le sont,
avec $u_{i+3} = \alpha_i u_i$ et $u_{i+2} = -\alpha_i u_i - u_{i+1}$.
Comparant les deux expressions de $u_{i+3}$ dans la base $u_i, u_{i+1}$,
\[
  u_{i+3} = -\alpha_{i+1} u_{i+1} - u_{i+2}
  = \alpha_i u_i + (1 - \alpha_{i+1})\, u_{i+1}
  \quad \text{et} \quad u_{i+3} = \alpha_i u_i ,
\]
d'où $\alpha_{i+1} = 1$ pour tout $i$. Ainsi $u_{i+3} = u_i$ et
$u_{i+1} + u_{i+2} = -u_i$ : le polygone est un triangle parcouru plusieurs
fois, et $n \equiv 0 \pmod 3$, car si $3$ ne divisait pas $n$ tous les $u_i$
seraient égaux, donc alignés\footnote{La page 24 fait ce calcul avec des
$\alpha_i$ puis des $\beta_i$ ; la liste $u_5, \ldots, u_9$ qu'elle écrit ne
suit pas la règle $u_{i+3} = \alpha_i u_i$ inscrite au-dessus (les indices de
$u_7, u_8, u_9$ sont ceux de la page, non corrigés par la transcription), et
l'on ne la reproduit pas. Elle conclut « $\beta_1 = 1$, i.e.\ $u_1 = u_4$ »,
« donc $u_i = u_{i+3}$ », « donc $n \equiv 0$ (3) », le cas
« $u_i = u_{i+1} = \cdots\ \forall i$ » étant écarté. Le haut de la page 25
refait le cas du triangle : $u_1 + u_2 \parallel u_0$, $u_2 + u_0 \parallel u_1$,
$u_0 + u_1 \parallel u_2$ donnent $u_2 = -u_0 - u_1$.}.

\textbf{Proposition} (page 25). Soit $V$ un plan vectoriel sur $k$ et
$u_* = (u_i)_{i \in \mathbf{Z}/n\mathbf{Z}}$ une suite de vecteurs non nuls de
$V$, non tous portés par une même droite. Soit $N(u_*)$ l'espace des
$(\lambda_i) \in k^{\mathbf{Z}/n\mathbf{Z}}$ qui satisfont $(\ast)$. Considérons
\begin{enumerate}
  \item[(a)] il existe $(\lambda_i) \in N(u_*)$ dont tous les $\lambda_i$ sont
    non nuls, et $\dim_k N(u_*) \geqslant 2$ ;
  \item[(b)] pour tout $i$, $u_i$, $u_{i+1} + u_{i+2}$ et $u_{i+3}$ sont
    colinéaires ;
  \item[(c)] pour tout $i$, $u_{i+3} = u_i$ et $u_{i+1} + u_{i+2} = -u_i$ ;
  \item[(d)] $n \equiv 0 \pmod 3$, $u_i = u_{i+3}$ pour tout $i$, et
    $u_0 + u_1 + u_2 = 0$.
\end{enumerate}
Alors (b), (c) et (d) sont équivalentes, (a) les entraîne, et elles entraînent
(a) si $k$ a au moins trois éléments. Dans ce cas $\dim N(u_*) = 2$, et
$N(u_*)$ est formé des $(\lambda_i)$ de période 3 tels que
$\lambda_0 + \lambda_1 + \lambda_2 = 0$\footnote{La Proposition est réécrite
trois fois sur place, et ses couches biffées sont en grande partie illisibles ;
on la donne comme elle se lit en dernier, avec trois changements. (1)
L'hypothèse de non-alignement est ajoutée : la suite $(e, -e, e, -e)$, un
segment parcouru deux fois dans les deux sens, a $N(u_*)$ de dimension 2,
formé des $(a, b, a, b)$, avec des multiplicateurs tous non nuls ; elle
satisfait (a) et (b), mais ni (c) ni (d). (2) La condition $n \equiv 0$ (3) est
mise dans (d) : la page l'y écrivait et l'a biffée, et sans elle la période 3 des
$\lambda_i$ n'a pas de sens sur $\mathbf{Z}/n\mathbf{Z}$. (3) La page suppose
$k$ infini pour la démonstration ; on dit ce qui sert, et où. Les lettres (a)
à (d) sont cerclées.}.

\pagerange{26}{27}
\section*{La démonstration (pages 26 et 27)}

\emph{(a) entraîne (b).} D'après l'injectivité de
$(\lambda_i) \mapsto (\lambda_1, \lambda_2, \lambda_3)$, $N(u_*)$ s'identifie à
son image dans le noyau de
\[
  k^3 \longrightarrow V, \qquad (\lambda_1, \lambda_2, \lambda_3) \longmapsto
  \lambda_1 u_0 - \lambda_2 (u_1 + u_2) + \lambda_3 u_3 ,
\]
donc
\[
  \dim N(u_*) \leqslant 3 - \operatorname{rg}(u_0,\ u_1 + u_2,\ u_3) .
\]
Si $\dim N(u_*) \geqslant 2$, ce rang est au plus 1 ; il vaut 1 puisque
$u_0 \neq 0$. Par permutation circulaire, $u_i$, $u_{i+1} + u_{i+2}$, $u_{i+3}$
sont colinéaires pour tout $i$. Cette implication ne se sert ni de la
non-nullité des $\lambda_i$, ni de l'hypothèse de non-alignement\footnote{La
page établit d'abord l'injectivité de $N(u_*) \to k^3$ par un détour : un
élément de $N(u_*)$ à composantes toutes non nulles est déterminé par trois
composantes consécutives ; le noyau a donc des translatés dans $N(u_*)$ qui ne
rencontrent le complémentaire de la réunion des hyperplans
$H_i = \{\lambda_i = 0\}$ qu'en un point ; « le corps $k$ étant infini », le
noyau est nul. L'argument est correct sous l'hypothèse (a) et pour $k$ infini,
et c'est pour lui que la marge porte « OPS $k$ infini » ; la phrase n'est
reconstituée qu'au prix de plusieurs lectures douteuses. L'argument direct
donné plus haut (deux composantes suffisent dès que les $u_i$ sont non nuls)
le rend inutile. La page montre aussi que l'image ne contient pas
$(1, 0, 0)$, sans quoi $u_0 = 0$, d'où $\dim N(u_*) \neq 3$ ; c'est
désormais immédiat. Elle écrit « $\operatorname{rg} N(u_*) = 2$ » pour
$\dim N(u_*) = 2$, et « $\operatorname{rg}$ » en surcharge de « dim ».}.

\emph{(b), (c) et (d) sont équivalentes.} (b) entraîne (c) : c'est le calcul
de la section précédente, qui suppose les $u_i$ non tous alignés. (c) entraîne
(d) en faisant $i = 0$, et $3 \mid n$ comme on l'a vu ; (d) entraîne (c) parce
que $u_i + u_{i+1} + u_{i+2}$, de période 3 en $i$, est une permutation
circulaire de $u_0 + u_1 + u_2$ ; (c) entraîne (b) trivialement\footnote{« Un
calcul facile plus haut montre b) $\Leftrightarrow$ c) et on voit de suite que
ceci équivaut à d) » : trois mots sont lus avec doute.}.

\emph{(d) entraîne (a).} Pour $(\lambda_i)$ de période 3, la relation $(\ast)$
devient, puisque $u_i + u_{i+1} = -u_{i-1}$, $u_{i+2} = u_{i-1}$ et
$\lambda_{i+2} = \lambda_{i-1}$,
\[
  (\lambda_{i-1} + \lambda_i + \lambda_{i+1})\, u_{i-1} = 0 ,
\]
c'est-à-dire $\lambda_0 + \lambda_1 + \lambda_2 = 0$. Ces suites forment un
plan, contenu dans $N(u_*)$, qui est de dimension au plus 2 : c'est
$N(u_*)$. Il contient un élément à composantes toutes non nulles dès que $k$ a
au moins trois éléments --- $(1, 1, -2)$ en caractéristique différente de 2,
$(1, \omega, \omega + 1)$ avec $\omega \notin \{0, 1\}$ en caractéristique 2 ;
sur $\mathbf{F}_2$ il n'y en a pas\footnote{La page calcule le plan des
$(\lambda_i)$ de période 3 et s'exclame : « Ils forment un vectoriel de
dim~2 ! » Elle écrit ensuite « Dém.\ d) $\Rightarrow$ a) » et le reste de la
page est blanc ; la démonstration est de l'édition.}.

\pagerange{28}{28}
\section*{Le Théorème (page 28)}

\textbf{Théorème.} Soit $V$ un plan vectoriel sur un corps infini $k$,
$n \geqslant 3$, et $u_* = (u_i) \in V^{\mathbf{Z}/n\mathbf{Z}}$ une suite de
vecteurs non nuls. Soit
\[
  N(u_*) = \bigl\{ (\lambda_i) \in k^{\mathbf{Z}/n\mathbf{Z}} \bigm|
  \lambda_i u_{i-1} - \lambda_{i+1} (u_i + u_{i+1}) + \lambda_{i+2} u_{i+2} = 0
  \ \ \forall\, i \bigr\},
\]
et $H_i \subset k^{\mathbf{Z}/n\mathbf{Z}}$ l'hyperplan de coordonnées
$\lambda_i = 0$. On est dans un et un seul des cas suivants.
\begin{enumerate}
  \item[(a)] \emph{Cas général} : $N(u_*) = 0$. Si $\sum u_i = 0$, cela
    signifie que l'application $\varphi$, des polygones fermés vers les aires
    de leurs oreilles, est une submersion en $u_*$.
  \item[(b)] $N(u_*) \not\subset \bigcup_i H_i$ ; alors
    $\dim N(u_*) \in \{1, 2\}$, et l'on distingue :
  \item[(b$_1$)] \emph{Cas spécial I} : $\dim N(u_*) = 1$ ;
  \item[(b$_2$)] \emph{Cas hyperspécial} : $\dim N(u_*) = 2$, ce qui, si
    les $u_i$ ne sont pas tous alignés, équivaut à : $n \equiv 0
    \pmod 3$, $u_i = u_{i+3}$ pour tout $i$, et $u_0 + u_1 + u_2 = 0$.
  \item[(c)] \emph{Cas spécial II} : $N(u_*) \neq 0$ et il existe $i$ tel
    que $N(u_*) \subset H_i$ ; alors $\dim N(u_*) = 1$.
\end{enumerate}
Les noms des cas sont de sa main, dans la marge\footnote{Entre guillemets et
en oblique ; « Cas hyperspécial » est souligné. La page écrit d'abord
« $n$ entier $\geqslant 1$ », le 1 sur un autre chiffre, peut-être un 3 ; on
garde $n \geqslant 3$. La relation qui définit $N(u_*)$ est serrée et
surchargée, et ses indices sont lus d'après la page 25.}.

\emph{Démonstration.} Si $N(u_*) \neq 0$, ou bien $N(u_*)$ n'est pas contenu
dans la réunion des $H_i$, ou bien il l'est, et alors il est contenu dans l'un
d'eux : un espace vectoriel sur un corps infini n'est pas réunion finie de
sous-espaces stricts. C'est le « Lemme ($N(u_*) \not\subset \bigcup H_i$) » de
la marge de la page 26, et c'est là, et là seulement, que sert l'hypothèse
$k$ infini. On a toujours $\dim N(u_*) \leqslant 2$ (page 23). Si
$\dim N(u_*) = 2$, l'injection $N(u_*) \to k^2$,
$\lambda \mapsto (\lambda_i, \lambda_{i+1})$, est bijective pour tout $i$, et
$N(u_*)$ n'est contenu dans aucun $H_i$ : en (c), la dimension est 1. Enfin
(b$_2$) est la Proposition\footnote{Comme dans la Proposition, l'équivalence de
(b$_2$) avec la condition de droite demande que les $u_i$ ne soient pas tous
alignés ; la page ne le dit pas. La glose de (a) est reconstituée : la page
écrit « [i.e., si $\sum u_i = 0$, l'application can.\
$\varphi : X(V) \to (\det V)^{\mathbf{Z}/n\mathbf{Z}}$ \ldots\ lin.\ \ldots\
$u_*$] », avec plusieurs mots illisibles, « lin. » lu avec doute et la lettre
devant « $(V)$ » surchargée ; la lecture « submersion », qui est ce que dit la
section des pages 14 et 15, est la nôtre.}.

\emph{Ce qui se passe pour les polygones fermés.} Si $\sum u_i = 0$ et que les
$u_i$ engendrent $V$, le groupe $SL(V)$ agit sur les polygones fermés en
préservant les $\mu_i$ ; son orbite, de dimension 3, est dans une fibre de
$\varphi$, si bien que le rang de $\varphi$ en $u_*$ est au plus $2n - 5$ et
$\dim N(u_*) \geqslant 5 - n$. Un triangle est donc toujours hyperspécial, un
quadrilatère non aligné toujours dans le cas spécial I, avec le multiplicateur
$(1, -1, 1, -1)$ des pages 16 et 17, et le cas général ne se rencontre qu'à
partir du pentagone --- où il est, de fait, le cas d'un polygone pris au
hasard\footnote{Ce paragraphe est le nôtre ; la page n'en dit rien. Le
dernier point a été vérifié par le calcul sur des polygones fermés pris au
hasard, pour $n = 5, 6, 7$.}.

\emph{Le critère du cas spécial II.} Supposons qu'un $\lambda \in N(u_*)$ non
nul ait une composante nulle, et numérotons pour qu'elle soit $\lambda_0$.
Alors $\lambda_1 \neq 0$ (deux composantes consécutives nulles annulent tout),
et l'on normalise $\lambda_1 = 1$. Le vecteur constant de $(L)$ est
$w = \lambda_0 u_{-1} - \lambda_1 u_1 = -u_1$, et l'on a :

\textbf{Proposition.} Il existe $\lambda \in N(u_*)$ avec $\lambda_0 = 0$ et
$\lambda_1 = 1$ si et seulement si
\begin{enumerate}
  \item[a)] $\sum u_i = 0$, et
  \item[b)] pour $2 \leqslant i \leqslant n-1$, ou bien $s_i = s_0$, ou bien
    la paire de directions de $\{u_{i-1}, u_i\}$ est celle de
    $\{u_1, u_{0i}\}$.
\end{enumerate}
Les $\lambda_i$ sont alors définis par $u_1 \cdot u_{0i} = \lambda_i\, u_{i-1}
\cdot u_i$ dans $\operatorname{Sym}^2 V$, et les indices $i$ tels que
$\lambda_i = 0$ sont exactement ceux pour lesquels $s_i = s_0$.

En effet, notons $(S_i)$ l'identité $\lambda_i\, u_{i-1} \cdot u_i = u_1 \cdot
u_{0i}$. $(S_1)$ est vraie. Si $(S_i)$ est vraie, la relation $(L)$ en $i$,
$\lambda_{i+1} u_{i+1} = \lambda_i u_{i-1} + u_1$, multipliée par $u_i$, donne
$(S_{i+1})$ ; inversement, $(S_i)$ et $(S_{i+1})$ donnent $(L)$ en $i$ en
simplifiant par $u_i$, car $\operatorname{Sym} V$ est intègre. Ainsi $(L)$ pour
tout $i$ équivaut à $(S_2), \ldots, (S_n)$. Or $(S_n)$ s'écrit
$\lambda_0\, u_{n-1} \cdot u_0 = 0 = u_1 \cdot \sum_i u_i$, c'est-à-dire a) ; et
$(S_i)$, pour $2 \leqslant i \leqslant n-1$, a une solution $\lambda_i$ si et
seulement si b) a lieu, puisque dans un plan deux produits de vecteurs non nuls
sont proportionnels exactement quand ils portent la même paire de droites.
Enfin $\lambda_i = 0$ si et seulement si $u_1 \cdot u_{0i} = 0$, si et
seulement si $u_{0i} = 0$\footnote{La page écrit le critère sous la forme
« a) $\sum u_i = 0$ ; b) $\forall\, 1 \leqslant i \leqslant n-1$, on ait
$u_{i-1}, u_i \parallel u_1, u_{0i}$ », et « si on normalise ($\lambda_1 = 1$),
on aura $\lambda_i$ défini par $u_1 u_{0i} = \lambda_i u_i u_{i-1}$ » ; le
« il faut et il suffit » est lu avec doute, et la démonstration est de
l'édition. Le cas $i = 1$ de b) est vide. La phrase sur les zéros,
« les $i$ tels que $\lambda_i = 0$ sont ceux pour lesquels $s_i = s_{i_0}$
(i.e.\ $u_{i_0} + \cdots + u_{i-1} = 0$) », est sur la page. Le critère
caractérise l'existence d'un multiplicateur nul en $0$ : on est alors dans le
cas (c), ou dans le cas hyperspécial, où tous les $H_i$ coupent $N(u_*)$
suivant une droite.}.

L'ennéagone de la page 18 en est l'exemple : ses côtés d'indice impair ont la
direction de $u_1$, ceux d'indice pair celle de $s_i - s_0$, et aucun sommet
autre que $s_0$ ne revient en $s_0$. Une note de marge, reliée au critère par
une double flèche, en donnait une forme géométrique --- un point $o$, une
direction $\Delta$, des côtés parallèles à $\Delta$ ou passant par $o$ ---,
mais elle est en grande partie biffée et la fin en est
illisible\footnote{On y lit « $\exists\, o \in E$ », « $\Delta$ »,
« $u_{i-1} \parallel \Delta$ », « $u_i \parallel$ », et une numérotation
1°), 2°), 3°) lue avec doute. Ce qu'on en dit est une lecture d'ensemble, que
le critère démontré confirme pour l'exemple de la page 18, où $o = s_0$ et
$\Delta$ est la direction de $u_1$.}.

\end{document}
